\section{A pre-registered confirmation}
\label{sec:confirm}

The newest evidence in Sections~\ref{sec:ar1-levels} and~\ref{sec:software}---arch's studentized SPA test after its correction, the author's corrected implementation, and the best-of-\nz{trials} tests at small autocorrelations---came from checks run once without pre-registration, and the joint tests had been run at only one autocorrelation, 0.2. Run v2c repeats those checks on fresh seeds under a pre-registration, extends the joint tests to small autocorrelations, and tests Proposition~2(b). The pre-registration (signed commit \nolinkurl{\nz{conf.prereg.commit}}, \nz{conf.prereg.utc}) fixed \nz{conf.count} hypotheses with numeric bars, the scripts that run and score them, the seeds and the closed-form predictions, before any search was drawn. It has three parts, all with \nz{trials} strategies, \nz{obs} days and a nominal 5\% level:
\begin{description}
  \item[S, the nine families.] On \nz{conf.searches.py} new null searches per family, arch at the merge commit of pull request 882 with its defaults, at block lengths 8 and 22, beside this paper's fixed-variance SPA test and Reality Check on the same searches (Python); on \nz{conf.searches.js} other searches per family, the author's implementation as v1 ran it and as corrected, with the same bootstrap draws (JavaScript).
  \item[L, small autocorrelations.] Independent Gaussian AR(1) strategies with $\phi=0.02$, 0.05 and 0.1: the same joint tests on the same numbers of searches, and the four single-series tests of Table~\ref{tab:ar1-levels} on \nz{conf.searches.single} searches per level.
  \item[C, correlated strategies.] Strategies built as in Proposition~2(b) with $\rho=0.3$ and 0.6 and $\phi=0$, 0.05 and 0.1, \nz{conf.searches.corr} searches per cell, scored against equation~\eqref{eq:ar1corr}.
\end{description}
Both parts finished within the registered time limit (\nz{conf.minutes.py} and \nz{conf.minutes.js} minutes with twelve workers), so the shortening rule did not bind. Between the commit and the run both scripts were smoke-tested on another seed, as the pre-registration allows; no code changed. Because every seed is fixed in the committed pre-registration and the run is deterministic, a re-run reproduces the result files---the JavaScript part, run again with \nz{conf.rerun.workers} workers instead of \nz{conf.workers}, gave an identical result file apart from its run time---so a different outcome could not have been obtained by running again without changing the committed seeds. Of the \nz{conf.count} hypotheses, \nz{conf.held_count} held and one, H10, did not (Table~\ref{tab:predictions-confirm}).

\paragraph{Re-studentizing, confirmed.} Across the nine families arch's corrected test rejected \nz{conf.range.S.arch882_b8_consistent}\% of null searches at block length 8 (H1) and \nz{conf.range.S.arch882_b22_consistent}\% at its default of 22 (H2), against \nz{conf.range.S.fixed_spa_c_b8}\% and \nz{conf.range.S.fixed_spa_c_b22}\% for the fixed-variance test on the same searches; at block length 22 the fixed-variance test rejected at least \nz{conf.mingap.fixed_arch_b22} points more in every family (H3; Table~\ref{tab:confirm-joint}). Under negative skew, the family with the largest excess, the fixed-variance test rejected \nz{conf.size.fixed_spa_c_b8.skew_negative}\% and arch \nz{conf.size.arch882_b8_consistent.skew_negative}\% at block length 8, and the author's implementation went from \nz{conf.size.v1.spa_consistent.skew_negative}\% to \nz{conf.size.corrected.spa_consistent.skew_negative}\% (H6). The corrected implementation rejected \nz{conf.range.S.corrected_nonar}\% outside the AR(1) family and \nz{conf.size.corrected.spa_consistent.ar1}\% within it, never more than \nz{conf.maxdiff.corrected_rc} points from the Reality Check on the same searches (H4, H5), and its StepM rejected at the same rate as its SPA test in every design (H7). Not every corrected test sat inside the 99\% band (\nz{band99.3000}\% for \nz{conf.searches.py} searches): arch at block length 22 was above it in \nz{conf.above.arch882_b22_consistent}, the Reality Check at block length 22 in \nz{conf.above.rc_b22}, and the corrected implementation in \nz{conf.above.corrected.spa_consistent}. These residual excesses are small next to the fixed-variance test's and of the size the unstudentized Reality Check shows with long blocks and correlated strategies (Section~\ref{sec:snooping}); the corrected tests are close to their level, not exact.

\paragraph{Joint tests at small autocorrelations.} At $\phi=0.02$--0.1 the tests that resample whole blocks and re-studentize, or do not studentize, stayed close to their level, while the single-series tests did not (H11--H13): arch rejected \nz{conf.range.L.arch882_b8_consistent}\% at block length 8 and \nz{conf.range.L.arch882_b22_consistent}\% at 22, the Reality Check \nz{conf.range.L.rc_b8}\% and the corrected SPA test and StepM \nz{conf.range.L.corrected.spa_consistent}\%, against \nz{conf.L.haircut_bonferroni_one_sided.0.05}\% and \nz{conf.L.haircut_bonferroni_one_sided.0.1}\% for the one-sided haircut at $\phi=0.05$ and 0.1. The fixed-variance test was oversized at every level (\nz{conf.range.L.fixed_spa_c_b8}\% at block length 8, \nz{conf.range.L.fixed_spa_c_b22}\% at 22; v1's implementation \nz{conf.range.L.v1.spa_consistent}\%). A block bootstrap carries a small autocorrelation into its resamples; a $t$-statistic read against Student's $t$ does not.

\paragraph{Proposition 2 and Lo's factor.} The single-series sizes confirmed the earlier run: LET rejected \nz{conf.L.luck_trials.0.02}\%, \nz{conf.L.luck_trials.0.05}\% and \nz{conf.L.luck_trials.0.1}\% at $\phi=0.02$, 0.05 and 0.1 (H8), within \nz{conf.maxz.H9} Monte Carlo standard errors of Proposition~2(a) for both tests (H9; Table~\ref{tab:confirm-single}). H10 failed. It required Lo's factor to keep both tests inside the 99\% band at every level; the haircut version did (\nz{conf.Lrange.haircut_bonferroni_lo_one_sided}\%), but LET-Lo rejected \nz{conf.L.luck_trials_lo.0.05}\% at $\phi=0.05$, just above the band's upper edge, and \nz{conf.L.luck_trials_lo.0.02}\% at $\phi=0.02$. With the earlier run, which put it at the edge at the same two levels, the evidence is that Lo's factor at an estimated autocorrelation leaves LET about half a point liberal at small $\phi$---it removes most of the excess, not all of it.

\paragraph{Correlated strategies.} Equation~\eqref{eq:ar1corr} matched the simulated sizes within \nz{conf.maxz.H14} standard errors in all six cells (H14), and the independence formula overstated them at $\rho=0.6$ by at least \nz{conf.minz.H15} standard errors (H15). At $\rho=0.6$ the agreement has a pattern: the closed form was \nz{conf.z.rho06.phi0} standard errors above the simulation at $\phi=0$ and \nz{conf.z.rho06.phi01} below it at $\phi=0.1$, so the simulated size rose somewhat faster with $\phi$ than the Gaussian-copula approximation predicts. Correlation makes the Bonferroni-type tests conservative, but it does not protect them: at $\rho=0.6$, LET rejected \nz{conf.C.luck_trials.0.6.0}\% without autocorrelation and \nz{conf.C.luck_trials.0.6.0.1}\% at $\phi=0.1$, and at $\rho=0.3$ it rose from \nz{conf.C.luck_trials.0.3.0}\% to \nz{conf.C.luck_trials.0.3.0.1}\%. A small autocorrelation first spends the margin that correlation gives, then exceeds the level.
